\section{Evaluation Settings for Model and Loop Configurations}\label{app:exp-settings}

\paragraph{Models and provider routes.}
RQ1 (E1) isolates model choice by comparing eight model checkpoints under a fixed Claude Code loop implementation: Opus-4.7, GPT-5.5, GLM-5.1, DeepSeek-V4P, Gemini-3.1-Pro, Qwen3.6-Plus, Kimi-2.6, and Grok-4.1-FR.
RQ1 (E2) isolates loop implementation by fixing the model to \texttt{azure/gpt-5.4-20260305} and varying Claude Code, Codex, GitHub Copilot, OpenHands, SWE-agent, and mini-SWE-agent.
RQ1 (E3) examines within family model changes while pairing each family with its corresponding loop implementation for GPT, Claude Opus, and Qwen.
Provider access is normalized through routes implemented in \texttt{long\_horizon\_bench/agents/}, including an Azure OpenAI deployment with AD token auth (\texttt{api-version} \texttt{2025-04-01-preview}), the GitHub Copilot token exchange endpoint, a DashScope Anthropic compatible proxy, and direct Anthropic SDK access for Anthropic native models. This normalization separates provider transport from the loop engineering factors attributed to model choice and loop implementation.

\paragraph{Decoding and evaluation input.}
Each loop implementation inherits provider defaults unless its evaluation adapter records an explicit override. For the Anthropic SDK path used by Claude Code, we set \texttt{max\_tokens=128{,}000} and leave temperature and top p at their SDK defaults. The SWE-agent LiteLLM call uses \texttt{temperature=0}. Top p is disabled because Anthropic rejects requests that set both fields. mini-SWE-agent forwards \texttt{model.model\_kwargs.reasoning\_effort=high} to providers that accept it. Loop implementations receive the evaluation contract in Appendix~\ref{app:prompts}; no additional custom system prompt is supplied unless the evaluation adapter defines one.

\paragraph{Runtime budget and concurrency.}
The evaluation runtime imposes a per task wall clock budget $T_{\max}$ chosen per task ($\geq$2~h, capped at 24~h for the longest PR Sequence and Research Evolution tasks) and passes it to each loop implementation as \texttt{timeout\_sec}.
For loop implementations with their own iteration cap, we set the cap above the wall clock budget, making wall clock time the evaluation horizon. Claude Code SDK uses \texttt{max\_iterations=10{,}000}, and mini-SWE-agent uses \texttt{--cost-limit 0} (uncapped) with \texttt{--exit-immediately}.
Each bash tool call inside the container is independently capped at 600~s, separating local command stalls from sustained loop progress.
Trials run concurrently in a \texttt{ThreadPoolExecutor}. Per trial credentials are passed as explicit kwargs through the task runner, keeping cross trial authorization races outside attribution to the evaluated loop implementation.
